\section{Theory}\label{sec:theory}

One principle runs through the results on operations. A library of calls fixes what can be reported, its coverage; a rule applied to the library decides what is reported, its selection. On a fixed library, an operation lowers the loss exactly when, on the records or items it changes, the errors it repairs outweigh the correct outputs it breaks, at the exchange rate the loss sets. Section~\ref{sec:theory:identities} states the exact relations for record sets: deleting records, averaging their coordinates, and reading a different field of one response. Section~\ref{sec:theory:behaviour} gives a behavioural model of how models report objects across repeated calls, views and models, which says when a partner adds coverage; unlike the relations, it is an assumption to be checked. Section~\ref{sec:theory:select} bounds what selection over a fixed library can reach. Section~\ref{sec:theory:points} treats a single answer chosen among several calls: the three places where a miss can arise, the coverage an added source brings, and the net effect of switching an answer, which is the principle's form for clicks. Section~\ref{sec:theory:cert} states what outputs alone cannot certify. Table~\ref{tab:theory_use} lists, for each result the method uses, the quantity it measures, the decision it makes and where it is checked. Proofs and supporting results are in Appendix~\ref{app:proofs}.

\begin{table*}[t]
\centering\footnotesize\setlength{\tabcolsep}{3.5pt}
\caption{What each result is used for: the quantity it measures, the decision it makes in the method, and the check in Section~\ref{sec:results}.}\label{tab:theory_use}
\begin{tabular}{@{}>{\raggedright\arraybackslash}p{2.6cm}>{\raggedright\arraybackslash}p{3.2cm}>{\raggedright\arraybackslash}p{4.5cm}>{\raggedright\arraybackslash}p{4.9cm}@{}}\toprule
Result & Quantity & Decision & Check \\\midrule
Intersection threshold (Prop.~\ref{thm:intersect}) & marginal matching rate $p_1$ of unconfirmed records & confirm and delete for set error on over-reporting sources, keep a single call on under-reporting sources; label rule of Step~7 & from 5 labelled images the threshold predicts the sign on the other images in 83\% of draws; the Step~7 rule lowers the loss against the better decision from 0.43 (family default alone) to 0.17 errors per image (\S\ref{sec:res:set}) \\
Geometry term (Prop.~\ref{prop:geometry}) & selection and geometry terms & coordinate averaging is a separate switch: on by default, the deletion-only variant in the menu & averaging helps in 37 of 54 pairs and hurts in 13; the threshold's only two failures on averaged intersections are geometry reversals (\S\ref{sec:res:set}) \\
Two readouts (Prop.~\ref{prop:readout}) & gap between stated count and list count, without labels & bounds what switching readouts can gain, without labels & on aerial and document images, for GPT-6.1 and Sonnet 5.5, the gap is 0.1 to 7.4 against errors of 5 to 46, and switching readouts changes little (\S\ref{sec:res:readout}) \\
Partner (Prop.~\ref{prop:partner}) & confirmation rates $h_+$, $h_-$ & a partner must reject more false than true records; cross-model partners in the menu & cross-model partners separate the two better in 32 of 36 comparisons (\S\ref{sec:res:set}) \\
Coverage (Prop.~\ref{prop:coverage}) & rescue rate $\phi$ on the menu's misses & three different models, not three calls of one & a GPT-6.1 repeat rescues 18\%, then 7.5\%; Sonnet 5.5 29\%, Gemini 3.8 37\%; given GPT-6.1's 82 misses, rates from the 260 GUI items give 46 items with no correct candidate among the 600, against 43 observed (\S\ref{sec:res:choose}) \\
Three sources (Prop.~\ref{prop:three}) & $F$, $L$, $W$ & the steps of Algorithm~\ref{alg:clicks}; three sources, not four & a fourth source lowers $F$ from 43 to 36 and raises $L$ from 3 to 7 and $W$ from 10 to 11 (\S\ref{sec:res:choose}) \\
Switch (Prop.~\ref{prop:switch}) & $e_T$, $r_{\mathrm{fix}}$, $r_{\mathrm{brk}}$ on the items acted on & compare clusters only on alarmed items; no switch to the crop-first point & the score alone misses 5.8\% of single-cluster and 31.7\% of multi-cluster items; switching can gain at most 10 items (\S\ref{sec:res:choose}) \\
Alarm and no certificate (Props.~\ref{prop:alarm}, \ref{prop:nocert}) & disagreement between reports & flag errors without labels; choose configurations with labels of other sources & alarms evaluated by risk and acceptance coverage (\S\ref{sec:res:persist}); configurations by leaving one source out (\S\ref{sec:res:transfer}) \\
Risk ordering (Prop.~\ref{prop:order}) & gap between a selector and the constrained oracle & selectors fitted with labels are optional (Step~7), not a default & median gap 26.1 points under cross-fitting (\S\ref{sec:res:select}) \\
\bottomrule
\end{tabular}
\end{table*}

\subsection{Exact relations between operations and losses}\label{sec:theory:identities}

\begin{proposition}[Deletion and addition]\label{prop:identity}
Let $S\subseteq A$, let $\Delta c=c(A)-c(S)$ be the net decrease in scored records and $\Delta r=r(A)-r(S)$ the number of matches lost. Then exactly
\[
E_w(A)-E_w(S)=\cfp\,\Delta c-(\cfp+\cfn)\,\Delta r,
\]
and in particular $E(A)-E(S)=\Delta c-2\Delta r$. Addition is the same relation read in reverse.
\end{proposition}

\begin{proposition}[Intersection threshold]\label{thm:intersect}
Deleting $A\setminus S$ lowers weighted set error if and only if $\cfp\,\Delta c>(\cfp+\cfn)\,\Delta r$, that is, for $\Delta c>0$, iff $\Delta r/\Delta c<\cfp/(\cfp+\cfn)$. Adding a non-empty set lowers weighted set error iff $(\cfp+\cfn)$ times the matches gained exceeds $\cfp$ times the scored records added.
\end{proposition}

In words: deleting records helps when few enough of them were matched. Over a population of images with $\mathbb{E}[\Delta c]>0$, the expected change has the sign of $\cfp/(\cfp+\cfn)-p_1$, where $p_1:=\mathbb{E}[\Delta r]/\mathbb{E}[\Delta c]$ is the ratio of expected matches lost to expected records removed. We call $p_1$ the \emph{marginal matching rate} of the deleted records; it is not the average of per-image precisions, and it equals the precision of the deleted records only when no remaining record can take over a deleted record's match. For $E$ the threshold is $1/2$. This is the cost-sensitive decision threshold \citep{elkan2001costsensitive} applied to the records a partner does not confirm; what is specific to record sets is that $\Delta r$ can be smaller than the number of matched records deleted, so the relevant rate is $p_1$ and not their precision. The proposition does not assume that the confirming report comes from the same model, the same view, or any probability model. It concerns deletion only: the retained records keep their coordinates.

For F1 the same argument gives a per-image threshold of $\mathrm{F1}(A)/2$ in place of $1/2$ (Corollary~\ref{cor:f1} in Appendix~\ref{app:proofs}), the marginal form of the known result that the F1-optimal threshold of a calibrated classifier is half its optimal F1 \citep{lipton2014f1}. Because $\mathrm{F1}(A)/2\le1/2$, a deletion that raises F1 on an image also lowers $E$ there, but not conversely. This agrees with the defaults of Section~\ref{sec:res:transfer}: confirmation for set error on over-reporting sources, a single call for F1. In the deletion case, the threshold is the same exchange rate as the price of a false output in the companion analysis of decomposition \citep{shi2026zoomin}.

\begin{proposition}[Geometry term]\label{prop:geometry}
Let $S_0\subseteq A$ be the confirmed records with their original coordinates and $S_f$ the same clusters with averaged coordinates. Then
\begin{align*}
E(A)-E(S_f)=\ &\underbrace{[E(A)-E(S_0)]}_{\text{selection (Prop.~\ref{prop:identity})}}\\
&+\underbrace{[E(S_0)-E(S_f)]}_{\text{geometry}},
\end{align*}
and the geometry term equals $[c(S_0)-c(S_f)]-2[r(S_0)-r(S_f)]$.
\end{proposition}
By convexity, the mean of a cluster's members is no farther, in mean squared distance, from any fixed target than the members are on average. That does not make every averaged record closer than its best member, and matching uses a threshold, so the geometry term can have either sign. On DocLayNet with GPT-6.1 confirmed by Sonnet 5.5 (60 images), $E(A)=8.567$, $E(S_0)=7.767$ and $E(S_f)=8.867$: selection gains $0.800$, geometry loses $1.100$, and the intersection with averaged coordinates is worse than the single call although deletion alone is better.

\begin{proposition}[Two readouts of one response]\label{prop:readout}
Let a decoded response have delivered records $Z$ and stated counts $d\neq\varnothing$, and let $\mathrm{lc}(Z)_c=\#\{z\in Z:\mathrm{class}(z)=c\}$ be the list count. Both readouts come from the same response, so the difference between their errors $\|d-n\|_1$ and $\|\mathrm{lc}(Z)-n\|_1$ isolates the readout and holds generation fixed. The gap $\|d-\mathrm{lc}(Z)\|_1$ is observable without labels, and
\[
\max\{\|d-n\|_1,\|\mathrm{lc}(Z)-n\|_1\}\ \ge\ \tfrac12\|d-\mathrm{lc}(Z)\|_1,
\]
\[
\big|\,\|d-n\|_1-\|\mathrm{lc}(Z)-n\|_1\big|\ \le\ \|d-\mathrm{lc}(Z)\|_1 .
\]
\end{proposition}
A positive gap shows that at least one readout is wrong; it does not show which. A small gap shows that no choice between the two readouts can change the error by more than the gap, whatever the truth: choosing a readout can never gain more than the gap, and errors that both readouts share, such as objects missing from the list and from the statement alike, are beyond its reach. With a list count of 9 and a stated count of 10, the gap is 1 whether the truth is 9 or 10. Which readout to use is therefore an empirical question per model and task family; the method's rule (list count unless the list is degenerate) is supported by the measurements of Section~\ref{sec:res:readout}, not by this proposition.

\subsection{Behaviour under repetition, changed views and changed models}\label{sec:theory:behaviour}

\paragraph{Joint reporting law.} Fix a complete interface $\eta$: model, prompt, schema, setting, cap and view. A call with interface $\eta$ reports a true object $o$ with probability $q_{o,\eta}$, and a candidate false object $f$ (a location where false records tend to appear) with probability $q_{f,\eta}$. Given these propensities, repeated calls with the same $\eta$ are independent. This joint law, with $q$ free to depend on every part of $\eta$, is the behavioural model; the consequences below use only it, together with conditional independence given the propensities for calls with different interfaces wherever that is stated, and not the explanatory logit parametrization (object, model and view terms) of Appendix~\ref{app:param}. A reported record contributes to the loss only through the matching: emitting a record near $o$, producing a match, and adding a match that no other record already provides are different events.

Three consequences are used in the paper.
\emph{Persistence}: for fixed $\eta$, the probability of missing $o$ in all $K$ calls is $(1-q_{o,\eta})^K$, and averaged over objects it is $\mathbb{E}[(1-q_{o,\eta})^K]$, which tends to the share of objects with $q=0$. The decay is slow when $q$ is small for some objects, so at a fixed budget many misses repeat.
\emph{Diminishing coverage}: for calls that are exchangeable given the image, the expected number of targets matched by the union of $K$ calls is concave in $K$; this follows from the submodularity of coverage \citep{nemhauser1978submodular} and of matching rank (Lemma~\ref{lem:rank}). The statement concerns raw coverage, not the loss of a constructed set: repeated calls can still remove false records that do not repeat. For example, with 20 true records that are always reported and 80 candidate false records reported independently with probability 0.2, one call has expected $E=16$ and three-call unanimity has expected $E=0.64$.
\emph{Confirmation rates}: if the records of two calls refer to the same candidate objects and the two calls are conditionally independent given the propensities, a partner call $B$ confirms a record of call $A$ with probability
\[
h_\pm=\frac{\mathbb{E}[q_Aq_B\mid\text{true / false}]}{\mathbb{E}[q_A\mid\text{true / false}]},
\]
with $h_+$ for true records and $h_-$ for false ones. For a second call with the same interface $q_B=q_A$, so $h=\mathbb{E}[q_A^2]/\mathbb{E}[q_A]$. A review call that receives the first answer has a different interface and is not covered by this formula.

\begin{proposition}[When a different partner is better]\label{prop:partner}
Assume deletion-only confirmation, no match takeover (so $\Delta r$ equals the number of true records deleted), no ignore region or all of $A$'s false records outside it, and expected counts $n_+$ (true) and $n_-$ (false) among $A$'s records, with $h^B_\pm$ averaged over these records. Then the expected reduction in set error $E$ (equal costs) from replacing $A$ by $A\cap B$ is
\[
G(B)=n_-\,(1-h_-^B)-n_+\,(1-h_+^B),
\]
and a partner $B$ beats a same-interface partner $A'$ iff $n_-\,(h_-^{A'}-h_-^{B})>n_+\,(h_+^{A'}-h_+^{B})$.
\end{proposition}
The partner must reject more of $A$'s false records than it loses of $A$'s true ones, weighted by their counts. Neither similar overall quality nor a larger separation $h_+-h_-$ is necessary or sufficient: with $n_+=100$ and $n_-=10$, the pair $(h_+,h_-)=(0.95,0.85)$ gives $G=-3.5$, while the more separated pair $(0.5,0)$ gives $G=-40$; and a weaker partner whose errors avoid $A$'s false records can satisfy the inequality. Without the assumptions, the exact relation of Proposition~\ref{prop:identity} replaces these expectations.

\paragraph{Reports from several views.} Nothing in the chain assumes that a report comes from one call on the whole image. A view is part of the interface $\eta$, so a view planner (tiling, the safe decomposition of the companion paper, zoom) produces calls like any other, and their records, mapped to image coordinates, enter the same constructions and losses. Merging the reports of several views is itself a construction: the union of their records, with duplicates in the overlaps removed by the association rule. Confirmation across views is covered by Proposition~\ref{thm:intersect} and Proposition~\ref{prop:partner} with one change. A view can report only objects that it contains, so for a record lying wholly inside no partner view, $h_+^B=h_-^B=0$: deleting it is a decision about that record's own $p_1$, on which the partner carries no information. A view should therefore confirm or reject only the records it contains. A different view changes the propensities, but it need not share fewer persistent misses: the coverage it adds splits into a gain in how often objects are reported and a gain in complementarity, as for the rescue rate after Proposition~\ref{prop:coverage} (worked example in Appendix~\ref{app:param}). Whether a view helps as a partner is again the inequality of Proposition~\ref{prop:partner}, which also counts its false records.

\begin{proposition}[Alarm from disagreement]\label{prop:alarm}
Consider a point task in which a record matches a target of the same class within distance $\theta$, with no ignore region, and let $E$ be its set error. Let $U_{2\theta}(A,B)=|A|+|B|-2\,m_{2\theta}(A,B)$, where $m_{2\theta}$ is the matching size between two reports when records of the same class are paired within $2\theta$. Then
\[
E(A)+E(B)\ \ge\ U_{2\theta}(A,B),
\]
hence $\max\{E(A),E(B)\}\ge\tfrac12\,U_{2\theta}(A,B)$.
\end{proposition}
The bound needs no labels: when $U>0$, at least one of the two reports is wrong, although the bound does not say which. Constructed instances attain it, and it can be loose: two identical wrong reports give $U=0$. The click alarm in our experiments (distance above a share of the longer image side) is an empirical alarm for the hit-in-box loss, evaluated by risk-coverage curves.

\subsection{Construction and selection}\label{sec:theory:select}

Fix a feasible family $\SA$ and a loss. Define the oracle risk $R_{\mathrm{or}}=\mathbb{E}[\min_{S\in\SA}\ell(S,Y)]$, which uses the truth to pick the best set; the best risk $R^*_\Phi=\inf_g\mathbb{E}[\ell(g(\Phi),Y)]$ of a rule that maps features $\Phi$ computed from outputs alone into $\SA$; and the risk $R_{\hat g}$ of a fitted rule.

\begin{proposition}[Risk ordering]\label{prop:order}
$R_{\mathrm{or}}\le R^*_\Phi\le R_{\hat g}$, and
\begin{align*}
R_{\hat g}-R_{\mathrm{or}}=\ &\underbrace{(R^*_\Phi-R_{\mathrm{or}})}_{\text{information not in }\Phi}\\
&+\underbrace{(R_{\hat g}-R^*_\Phi)}_{\text{selector and estimation}}.
\end{align*}
\end{proposition}
An observed gap between a selector and the oracle bounds the sum but does not apportion it. A rule whose fitting uses labels is a \emph{label-fitted output-feature selector}: it reads no labels at inference, but its fitting does. The gap concerns a fixed library of calls; it says nothing about what additional calls or a different library would add.

For the family with at most one representative per cluster, which the selectors of Section~\ref{sec:res:select} use, the oracle is a maximum matching between clusters and targets and is computed exactly (Proposition~\ref{prop:oracle} in Appendix~\ref{app:proofs}); unrestricted matching over all records can exceed it but is not attainable within the family (Section~\ref{sec:res:select} reports both).

\subsection{Choosing one answer among several calls}\label{sec:theory:points}

A click has one answer, so the constructions above reduce to choosing one point. Let several calls (models, settings, views) give an item a menu $A$ of candidate points. Group the candidates into \emph{clusters}, the connected components of the graph that joins two candidates within distance $\delta$; fix one \emph{representative} per cluster before any choice is made; and let a rule output one representative. For clicks these clusters replace the association rule of Section~\ref{sec:setup:feasible}. The loss is the miss, $\mathbb{1}[\hat y\notin\text{target}]$. In the GUI-grounding tier of Section~\ref{sec:method:clicks} a scoring model, the \emph{scorer}, gives each candidate a hit score and compares representatives; Section~\ref{sec:res:choose} reports it on 600 click items.

\begin{proposition}[Three sources of a miss]\label{prop:three}
For every rule that outputs a representative, the miss rate is exactly $F+L+W$, where $F$ is the probability that no candidate in $A$ hits, $L$ that some candidate hits but no representative does, and $W$ that some representative hits but the rule outputs one that misses.
\end{proposition}

The terms locate where a miss arises; they are not independent effects of interventions. With the candidate menu fixed, changing the representatives leaves $F$ unchanged but can move a miss between $L$ and $W$. With candidates and representatives fixed, changing only the final choice acts on $W$. Only changing or enlarging the candidate menu can lower $F$, and doing so can change all three terms at once, because a new candidate can also change the clusters, the representatives and the comparison. Section~\ref{sec:res:choose} measures the three terms separately when a source is added. How much an added candidate lowers $F$ needs no independence assumption:

\begin{proposition}[Coverage of an added source]\label{prop:coverage}
Let $F$ be the probability that every candidate of a menu misses, and $\phi=\Pr(\text{an added candidate hits}\mid\text{every earlier candidate misses})$. With the added candidate, the probability becomes $F(1-\phi)$. If the candidates are conditionally independent given per-item hit probabilities $q_1,\dots,q_k$, then $F=\mathbb{E}[\prod_{j=1}^{k}(1-q_j)]$.
\end{proposition}

The relevant quantity is the rescue rate $\phi$ on the items the existing menu misses, not the added source's own accuracy or its correlation alone. For the indicators $I$ that some earlier candidate hits and $J$ that the added candidate hits, with $\Pr(I=0)>0$,
\[
\phi=\Pr(J=1\mid I=0)=\mathbb{E}[J]-\frac{\mathrm{Cov}(I,J)}{1-\mathbb{E}[I]},
\]
the binary form of the decomposition into reporting rate and complementarity for views (Appendix~\ref{app:param}). At equal hit rates, lower covariance raises the rescue rate; when hit rates differ, both terms matter. Persistence makes the covariance of a repeated call with the same interface large, which is why a repeat tends to rescue little, but a weak partner can rescue less than a repeat. Section~\ref{sec:res:choose} therefore measures $\phi$ for each tested source.

\begin{proposition}[Net effect of a switch]\label{prop:switch}
Let a base rule output $b$, and let an action replace $b$ by $g$ on a set $T$ of items decided from outputs. With $e_T=\Pr(b\text{ misses}\mid T)$, the repair rate $r_{\mathrm{fix}}=\Pr(g\text{ hits}\mid b\text{ misses},T)$ and the break rate $r_{\mathrm{brk}}=\Pr(g\text{ misses}\mid b\text{ hits},T)$, the miss rate falls by exactly
\[
\Pr(T)\,\big[\,e_T\,r_{\mathrm{fix}}-(1-e_T)\,r_{\mathrm{brk}}\,\big].
\]
If, within an event $U\supseteq T$, membership in $T$ is independent of whether $b$ misses, then $e_T=\Pr(b\text{ misses}\mid U)$.
\end{proposition}

This is the click counterpart of Proposition~\ref{thm:intersect}: an action pays when the errors it repairs outweigh the correct answers it breaks, counted on the items it touches. Disagreement alone does not say which candidate is wrong (Proposition~\ref{prop:alarm} bounds the sum of two errors, not the base error), so where to act is an empirical question. Where the base miss rate $e_T$ is low, as on items whose candidates agree, an action must break very few correct answers to pay: with $e_T=5.8\%$, fewer than about 3.1\% if it repairs half of the misses (Section~\ref{sec:res:choose} measures $e_T$ with and without an alarm). The scorer is a design choice of the configuration and not a consequence of the theory: choosing the best of $N_{\mathrm{s}}$ candidate scorers to within a miss-rate gap $\varepsilon_{\mathrm{s}}$ with probability $1-\beta_{\mathrm{s}}$ takes on the order of $\log(N_{\mathrm{s}}/\beta_{\mathrm{s}})/\varepsilon_{\mathrm{s}}^2$ labelled items under independent sampling, so the pipeline ships one scorer and does not choose it per task.

\subsection{What outputs alone can and cannot certify}\label{sec:theory:cert}

\begin{proposition}[No output-only certificate]\label{prop:nocert}
Fix a reported set $S$, a loss $L$ and a tolerance $\epsilon$. Suppose two admissible worlds induce the same law $P$ of the observable output history $H$, while $L(S,Y_0)\le\epsilon$ in the first and $L(S,Y_1)>\epsilon$ in the second. Then every output-only acceptance rule $a(H)\in[0,1]$, possibly randomized, accepts with the same probability in both worlds, so no such rule can accept with probability at least $1-\beta$ in the first world and at most $\alpha$ in the second if $\alpha+\beta<1$. More generally, $P_1(\text{accept})\ge P_0(\text{accept})-\mathrm{TV}(P_0,P_1)$, where $\mathrm{TV}$ is the total variation distance.
\end{proposition}
The same argument applies to any property of $Y$ on which two such worlds disagree, for example whether $L(S,Y)<L(S',Y)$ for two reports $S$ and $S'$, so which of two reports is more accurate cannot be certified from outputs alone either. The argument is the two-point method of minimax lower bounds \citep{tsybakov2009nonparametric}; methods that estimate accuracy from the agreement of several classifiers without labels \citep{platanios2014unlabeled} rely on assumptions, such as weakly dependent errors, that rule such pairs out. Such pairs arise from objects that a model rarely reports. If the bad world adds one target that each of $K$ calls reports with probability at most $\vartheta$, given the earlier output, and the target does not change the rest of the output, then $\mathrm{TV}(P_0,P_1)\le K\vartheta$. Any test's advantage in telling the two worlds apart, the difference between its acceptance probabilities in them, is then at most $K\vartheta$ (with equal priors, the best probability of naming the right world is $(1+\mathrm{TV})/2$), and no test is reliable uniformly over admissible classes in which $\vartheta$ can be arbitrarily small. The practical reading: agreement alone does not certify correctness. This does not prevent running a default procedure, ordering interfaces under explicit assumptions tested on other data, as the companion paper does for decomposition \citep{shi2026zoomin}, or operating a configuration chosen on source data without new labels.

\begin{corollary}[Sign of a count gap]\label{cor:nonid}
Suppose that, for a given report $S$, some admissible world has $|S|>n$, and that the admissible worlds are closed under adding any number of targets that no interface reports and that do not change the rest of the output. Then the sign of $|S|-n$ is not determined by the outputs: that world has the same output law as one in which $S$ under-counts. (This differs from the sign of a recall gain between two interfaces, which assumptions such as monotonicity in tokens per object can fix.)
\end{corollary}
Without further information, neither deleting unconfirmed records nor keeping them dominates for a count objective. Our count procedure keeps them, a conservative choice whose quality we evaluate empirically.
